\documentclass{article}
\usepackage{amsmath,amssymb,amsthm}
\usepackage{enumitem}
\usepackage{hyperref}
\usepackage{geometry}
\geometry{margin=1in}
\newtheorem{theorem}{Theorem}
\newtheorem{lemma}{Lemma}
\newtheorem{corollary}{Corollary}
\newcommand{\E}{\mathbb{E}}
\newcommand{\R}{\mathbb{R}}
\newcommand{\norm}[1]{\left\|#1\right\|}
\newcommand{\inner}[1]{\left\langle#1\right\rangle}
\newcommand{\grad}{\nabla}
\begin{document}

\section*{Algorithm Name}
\textbf{Stochastic Variance Reduced Gradient (SVRG) – Double Loop Formulation}

\section*{Mathematical Formulation}
Let the finite‑sum objective be 
\[
f(x)\;=\;\frac{1}{n}\sum_{i=1}^{n} f_i(x),\qquad x\in\R^{d},
\]
where each component $f_i$ is convex and $L$–smooth.  
Denote by $x^{*}$ the unique minimizer of $f$ (strong convexity $\mu>0$).

\medskip
\noindent\textbf{Outer loop (epoch $s=0,1,\dots$)}  
\[
\tilde{x}^{s}\;\leftarrow\;x\; \text{(snapshot point)} ,
\qquad 
\tilde{g}^{s}\;=\;\grad f(\tilde{x}^{s})\;=\;\frac1n\sum_{i=1}^{n}\grad f_i(\tilde{x}^{s}).
\]

\medskip
\noindent\textbf{Inner loop (inner iteration $t=0,\dots,m-1$)}  
Initialize $x_{0}^{s}=\tilde{x}^{s}$.  
For each $t$ draw an index $i_{t}^{s}\!\in\!\{1,\dots,n\}$ uniformly at random and form the
variance‑reduced estimator
\[
v_{t}^{s}\;=\;\grad f_{i_{t}^{s}}(x_{t}^{s})\;-\;\grad f_{i_{t}^{s}}(\tilde{x}^{s})\;+\;\tilde{g}^{s}.
\tag{1}
\]
Update
\[
x_{t+1}^{s}\;=\;x_{t}^{s}\;-\;\eta\,v_{t}^{s},
\tag{2}
\]
where $\eta>0$ is the stepsize.

After $m$ inner steps we set the next snapshot as the average (or a uniformly
random inner iterate):
\[
\tilde{x}^{s+1}\;=\;\frac{1}{m}\sum_{t=0}^{m-1}x_{t}^{s}
\quad\text{(or } \tilde{x}^{s+1}=x_{U}^{s}\text{ with }U\sim\text{Unif}\{0,\dots,m-1\}\text{)}.
\]

\section*{Pseudocode}
\begin{enumerate}[label=\textbf{Step \arabic*:}, leftmargin=*, itemsep=0.3ex]
\item \textbf{Input:} stepsize $\eta$, inner loop length $m$, number of epochs $S$,
initial point $x^{0}$.
\item \textbf{Initialize:} $\tilde{x}^{0}\gets x^{0}$.
\item \textbf{For} $s=0,\dots,S-1$ \textbf{do}
\begin{enumerate}[label=\arabic*.]
\item Compute the \emph{full} gradient at the snapshot:
\[
\tilde{g}^{s}\gets \frac{1}{n}\sum_{i=1}^{n}\grad f_i(\tilde{x}^{s}).
\]
\item Set $x_{0}^{s}\gets\tilde{x}^{s}$.
\item \textbf{For} $t=0,\dots,m-1$ \textbf{do}
\begin{enumerate}[label=\alph*.]
\item Sample $i_{t}^{s}\sim\text{Uniform}\{1,\dots,n\}$.
\item Compute the variance‑reduced estimator (1).
\item Update $x_{t+1}^{s}=x_{t}^{s}-\eta v_{t}^{s}$ (2).
\end{enumerate}
\item End inner loop.
\item Set the next snapshot, e.g.
\[
\tilde{x}^{s+1}\gets\frac{1}{m}\sum_{t=0}^{m-1}x_{t}^{s}.
\]
\end{enumerate}
\item End outer loop.
\item \textbf{Output:} $\tilde{x}^{S}$.
\end{enumerate}

\section*{Dry Run Example}
Consider a single‑sample problem $n=1$ with
\[
f(w)=\bigl(w-3\bigr)^{2},\qquad \grad f(w)=2\bigl(w-3\bigr).
\]
Take $w_{0}=0$, stepsize $\eta=0.1$, inner length $m=3$ (only one epoch is needed because
$n=1$).  

\medskip
\noindent\textbf{Epoch $0$ (snapshot $\tilde w^{0}=0$)}  
Full gradient: $\tilde g^{0}=2(0-3)=-6$.

\medskip
\noindent\textbf{Inner iteration $t=0$}
\[
v_{0}= \grad f_{1}(w_{0})-\grad f_{1}(\tilde w^{0})+\tilde g^{0}=(-6)-(-6)+(-6)=-6,
\]
\[
w_{1}=w_{0}-\eta v_{0}=0-0.1(-6)=0.6,
\qquad f(w_{1})=(0.6-3)^{2}=5.76.
\]

\medskip
\noindent\textbf{Inner iteration $t=1$}
\[
v_{1}= \grad f_{1}(w_{1})-\grad f_{1}(\tilde w^{0})+\tilde g^{0}
=2(0.6-3)-(-6)-6=-4.8-(-6)-6=-4.8+6-6=-4.8,
\]
\[
w_{2}=w_{1}-\eta v_{1}=0.6-0.1(-4.8)=1.08,
\qquad f(w_{2})=(1.08-3)^{2}=3.6864.
\]

\medskip
\noindent\textbf{Inner iteration $t=2$}
\[
v_{2}=2(1.08-3)-(-6)-6=-3.84+6-6=-3.84,
\]
\[
w_{3}=w_{2}-\eta v_{2}=1.08-0.1(-3.84)=1.464,
\qquad f(w_{3})=(1.464-3)^{2}=2.366.
\]

\medskip
\noindent\textbf{Snapshot update}  
\[
\tilde w^{1}= \frac{1}{3}(w_{0}+w_{1}+w_{2}) = \frac{0+0.6+1.08}{3}=0.56,
\qquad f(\tilde w^{1})\approx 5.9536.
\]
The iterates move steadily toward the optimum $w^{*}=3$.

\section*{Assumptions}
\begin{enumerate}[label=\textbf{A\arabic*:}, leftmargin=*, itemsep=0.3ex]
\item \textbf{Smoothness.} Each $f_i$ is $L$‑smooth:
\[
\forall x,y\in\R^{d},\qquad 
f_i(y)\le f_i(x)+\inner{\grad f_i(x)}{y-x}+ \frac{L}{2}\norm{y-x}^{2}.
\tag{3}
\]
\item \textbf{Strong Convexity.} $f$ is $\mu$‑strongly convex ($\mu>0$):
\[
\forall x,y\in\R^{d},\qquad 
f(y)\ge f(x)+\inner{\grad f(x)}{y-x}+ \frac{\mu}{2}\norm{y-x}^{2}.
\tag{4}
\]
\item \textbf{Finite Sum.} $f(x)=\frac1n\sum_{i=1}^{n}f_i(x)$ with $n\ge 1$.
\item \textbf{Stepsize.} $\eta$ satisfies $0<\eta<\frac{1}{2L}$.
\item \textbf{Inner Loop Length.} $m\ge 1$ (often $m=2n$ in practice).
\end{enumerate}

\section*{Main Convergence Theorem}
\begin{theorem}[Linear convergence of SVRG]\label{thm:linear}
Under Assumptions \textbf{A1}–\textbf{A5}, let the snapshot $\tilde{x}^{s+1}$ be the
average of the $m$ inner iterates. Define the contraction factor
\[
\rho \;:=\; \frac{1}{\mu\eta(1-2L\eta)}\;\frac{1}{m}
\;+\; \bigl(1-\tfrac1m\bigr)\bigl(1-2\mu\eta\bigr).
\tag{5}
\]
If $0<\eta\le \frac{1}{4L}$ and $m\ge \frac{2L}{\mu}$, then $0<\rho<1$ and
\[
\E\!\bigl[f(\tilde{x}^{s+1})-f(x^{*})\bigr]\;\le\;\rho\;
\E\!\bigl[f(\tilde{x}^{s})-f(x^{*})\bigr],\qquad\forall s\ge0.
\tag{6}
\]
Consequently,
\[
\E\!\bigl[f(\tilde{x}^{S})-f(x^{*})\bigr]\;\le\;\rho^{S}\,
\bigl[f(\tilde{x}^{0})-f(x^{*})\bigr].
\tag{7}
\]
\end{theorem}

\section*{Theoretical Properties}
\begin{enumerate}[label=\textbf{P\arabic*:}, leftmargin=*, itemsep=0.3ex]
\item \textbf{Stability.} The contraction factor $\rho$ is decreasing in $\eta$
until the bound $\eta<1/(2L)$, guaranteeing stable progress.
\item \textbf{Hyperparameter Sensitivity.} Choosing $\eta=1/(4L)$ and $m=2n$ yields
$\rho\le 3/4$, i.e. a constant‑factor reduction per epoch independent of $n$.
\item \textbf{Comparison.} Compared with plain SGD (sublinear $O(1/t)$ rate) SVRG attains a
geometric rate for strongly convex problems, matching the optimal rate of
full‑gradient methods while using only $n+m$ gradient evaluations per epoch.
\end{enumerate}

\section*{Discussion}
The theorem shows that the double‑loop SVRG algorithm enjoys a linear
convergence rate under standard smoothness and strong‑convexity assumptions.
The proof relies on the unbiasedness of the variance‑reduced estimator and a
tight bound on its second moment. The constants in $\rho$ highlight the
trade‑off between stepsize $\eta$ and inner‑loop length $m$.

\section*{Preliminaries (Definitions and Lemmas)}
\begin{lemma}[Unbiased estimator]\label{lem:unbiased}
For any $x$, snapshot $\tilde{x}$ and random index $i\sim\mathrm{Unif}\{1,\dots,n\}$,
\[
\E_{i}\!\bigl[v(x,\tilde{x})\mid x,\tilde{x}\bigr]
\;=\;\grad f(x).
\tag{8}
\]
\end{lemma}
\begin{proof}
\[
\E_{i}[v]=\frac1n\sum_{i=1}^{n}
\bigl(\grad f_i(x)-\grad f_i(\tilde{x})\bigr)+\grad f(\tilde{x})
\;=\;\grad f(x)-\grad f(\tilde{x})+\grad f(\tilde{x})
\;=\;\grad f(x).
\] 
\end{proof}

\begin{lemma}[Variance bound]\label{lem:var}
For any $x,\tilde{x}$,
\[
\E_{i}\!\bigl[\|v(x,\tilde{x})-\grad f(x)\|^{2}\mid x,\tilde{x}\bigr]
\;\le\;2L\bigl(f(x)-f(x^{*})+f(\tilde{x})-f(x^{*})\bigr).
\tag{9}
\]
\end{lemma}
\begin{proof}
Using $L$‑smoothness and the identity
\[
v-\grad f(x)=\bigl(\grad f_i(x)-\grad f_i(\tilde{x})\bigr)-\bigl(\grad f(x)-\grad f(\tilde{x})\bigr),
\]
we have
\[
\|v-\grad f(x)\|^{2}
\le 2\|\grad f_i(x)-\grad f_i(\tilde{x})\|^{2}
   +2\|\grad f(x)-\grad f(\tilde{x})\|^{2}.
\tag{10}
\]
By $L$‑smoothness,
\[
\|\grad f_i(x)-\grad f_i(\tilde{x})\|^{2}\le 2L\bigl(f_i(x)-f_i(\tilde{x})-\inner{\grad f_i(\tilde{x})}{x-\tilde{x}}\bigr)
\le 2L\bigl(f_i(x)-f_i(x^{*})\bigr)+2L\bigl(f_i(\tilde{x})-f_i(x^{*})\bigr).
\tag{11}
\]
Averaging over $i$ and using the same bound for the full gradient term gives (9). 
\end{proof}

\section*{Main Proof (Step‑by‑Step Derivation)}
We prove Theorem~\ref{thm:linear}.  Throughout the proof the expectation $\E$ is taken
over all random choices of indices $i_{t}^{s}$ and the random choice of the
snapshot (if the average is replaced by a uniform inner iterate).

\medskip
\noindent\textbf{Step 1 (One inner update).}  
From the update rule (2) and smoothness (3) we have for any $x^{*}$
\[
f(x_{t+1}^{s})\le f(x_{t}^{s})+\inner{\grad f(x_{t}^{s})}{x_{t+1}^{s}-x_{t}^{s}}
+\frac{L}{2}\norm{x_{t+1}^{s}-x_{t}^{s}}^{2}.
\tag{12}
\]
Since $x_{t+1}^{s}=x_{t}^{s}-\eta v_{t}^{s}$, (12) becomes
\[
f(x_{t+1}^{s})\le f(x_{t}^{s})-\eta\inner{\grad f(x_{t}^{s})}{v_{t}^{s}}
+\frac{L\eta^{2}}{2}\norm{v_{t}^{s}}^{2}.
\tag{13}
\]

\medskip
\noindent\textbf{Step 2 (Insert unbiasedness).}  
Take conditional expectation w.r.t. $i_{t}^{s}$ and use Lemma~\ref{lem:unbiased}:
\[
\E_{i_{t}^{s}}\!\bigl[\inner{\grad f(x_{t}^{s})}{v_{t}^{s}}\mid x_{t}^{s},\tilde{x}^{s}\bigr]
=\norm{\grad f(x_{t}^{s})}^{2}.
\tag{14}
\]

\medskip
\noindent\textbf{Step 3 (Second‑moment expansion).}  
We write
\[
\norm{v_{t}^{s}}^{2}
=\norm{\grad f(x_{t}^{s})}^{2}
+ \norm{v_{t}^{s}-\grad f(x_{t}^{s})}^{2}
+2\inner{\grad f(x_{t}^{s})}{v_{t}^{s}-\grad f(x_{t}^{s})}.
\tag{15}
\]
Taking expectation and using Lemma~\ref{lem:unbiased} again eliminates the cross term,
\[
\E_{i_{t}^{s}}\!\bigl[\norm{v_{t}^{s}}^{2}\mid x_{t}^{s},\tilde{x}^{s}\bigr]
=\norm{\grad f(x_{t}^{s})}^{2}
+\E_{i_{t}^{s}}\!\bigl[\norm{v_{t}^{s}-\grad f(x_{t}^{s})}^{2}\mid x_{t}^{s},\tilde{x}^{s}\bigr].
\tag{16}
\]

\medskip
\noindent\textbf{Step 4 (Apply variance bound).}  
Insert Lemma~\ref{lem:var} into (16):
\[
\E_{i_{t}^{s}}\!\bigl[\norm{v_{t}^{s}}^{2}\mid x_{t}^{s},\tilde{x}^{s}\bigr]
\le \norm{\grad f(x_{t}^{s})}^{2}
+2L\bigl(f(x_{t}^{s})-f(x^{*})+f(\tilde{x}^{s})-f(x^{*})\bigr).
\tag{17}
\]

\medskip
\noindent\textbf{Step 5 (Take full expectation).}  
Plug (14) and (17) into (13) and then take expectation over $i_{t}^{s}$:
\[
\begin{aligned}
\E\!\bigl[f(x_{t+1}^{s})\mid x_{t}^{s},\tilde{x}^{s}\bigr]
&\le f(x_{t}^{s})-\eta\norm{\grad f(x_{t}^{s})}^{2}
+\frac{L\eta^{2}}{2}\Bigl(\norm{\grad f(x_{t}^{s})}^{2}
+2L\bigl(f(x_{t}^{s})-f(x^{*})+f(\tilde{x}^{s})-f(x^{*})\bigr)\Bigr)\\[0.2em]
&= f(x_{t}^{s})-\Bigl(\eta-\frac{L\eta^{2}}{2}\Bigr)\norm{\grad f(x_{t}^{s})}^{2}
+L^{2}\eta^{2}\bigl(f(x_{t}^{s})-f(x^{*})+f(\tilde{x}^{s})-f(x^{*})\bigr).
\end{aligned}
\tag{18}
\]

\medskip
\noindent\textbf{Step 6 (Strong convexity to replace gradient norm).}  
From $\mu$‑strong convexity (4) we have
\[
\frac{1}{2L}\norm{\grad f(x_{t}^{s})}^{2}\;\ge\; f(x_{t}^{s})-f(x^{*})-\frac{L-\mu}{2L}\norm{x_{t}^{s}-x^{*}}^{2}.
\tag{19}
\]
A simpler inequality that suffices for the linear rate is
\[
\norm{\grad f(x_{t}^{s})}^{2}\;\ge\;2\mu\bigl(f(x_{t}^{s})-f(x^{*})\bigr).
\tag{20}
\]

\medskip
\noindent\textbf{Step 7 (Insert (20) into (18)).}  
Using (20) we obtain
\[
\begin{aligned}
\E\!\bigl[f(x_{t+1}^{s})\mid x_{t}^{s},\tilde{x}^{s}\bigr]
&\le f(x_{t}^{s})-\bigl(\eta- L\eta^{2}\bigr)\mu\bigl(f(x_{t}^{s})-f(x^{*})\bigr)\\
&\qquad+L^{2}\eta^{2}\bigl(f(x_{t}^{s})-f(x^{*})+f(\tilde{x}^{s})-f(x^{*})\bigr) .
\end{aligned}
\tag{21}
\]

\medskip
\noindent\textbf{Step 8 (Collect terms).}  
Define $\alpha:=\eta\mu(1- L\eta)-L^{2}\eta^{2}$. Since $0<\eta<\frac{1}{2L}$, $\alpha>0$.
Equation (21) rewrites as
\[
\E\!\bigl[f(x_{t+1}^{s})-f(x^{*})\mid x_{t}^{s},\tilde{x}^{s}\bigr]
\le\bigl(1-\alpha\bigr)\bigl(f(x_{t}^{s})-f(x^{*})\bigr)
+L^{2}\eta^{2}\bigl(f(\tilde{x}^{s})-f(x^{*})\bigr).
\tag{22}
\]

\medskip
\noindent\textbf{Step 9 (Recursive bound over the inner loop).}  
Unrolling (22) for $t=0,\dots,m-1$ and taking full expectation yields
\[
\E\!\bigl[f(x_{t}^{s})-f(x^{*})\bigr]
\le (1-\alpha)^{t}\bigl(f(\tilde{x}^{s})-f(x^{*})\bigr)
+L^{2}\eta^{2}\bigl(f(\tilde{x}^{s})-f(x^{*})\bigr)\sum_{k=0}^{t-1}(1-\alpha)^{k}.
\tag{23}
\]
The geometric sum evaluates to $\frac{1-(1-\alpha)^{t}}{\alpha}$, hence
\[
\E\!\bigl[f(x_{t}^{s})-f(x^{*})\bigr]
\le\Bigl[(1-\alpha)^{t}+ \frac{L^{2}\eta^{2}}{\alpha}\bigl(1-(1-\alpha)^{t}\bigr)\Bigr]
\bigl(f(\tilde{x}^{s})-f(x^{*})\bigr).
\tag{24}
\]

\medskip
\noindent\textbf{Step 10 (Average snapshot).}  
Recall $\tilde{x}^{s+1}= \frac{1}{m}\sum_{t=0}^{m-1}x_{t}^{s}$. By convexity of $f$,
\[
f(\tilde{x}^{s+1})\le \frac{1}{m}\sum_{t=0}^{m-1}f(x_{t}^{s}).
\tag{25}
\]
Taking expectation and using (24) with $t$ ranging from $0$ to $m-1$ gives
\[
\begin{aligned}
\E\!\bigl[f(\tilde{x}^{s+1})-f(x^{*})\bigr]
&\le \frac{1}{m}\sum_{t=0}^{m-1}
\Bigl[(1-\alpha)^{t}+ \frac{L^{2}\eta^{2}}{\alpha}\bigl(1-(1-\alpha)^{t}\bigr)\Bigr]\\
&\qquad\;\times \bigl(f(\tilde{x}^{s})-f(x^{*})\bigr).
\end{aligned}
\tag{26}
\]

\medskip
\noindent\textbf{Step 11 (Closed‑form of the coefficient).}  
The two sums are elementary:
\[
\frac{1}{m}\sum_{t=0}^{m-1}(1-\alpha)^{t}
= \frac{1-(1-\alpha)^{m}}{m\alpha},
\qquad
\frac{1}{m}\sum_{t=0}^{m-1}\bigl(1-(1-\alpha)^{t}\bigr)
=1-\frac{1-(1-\alpha)^{m}}{m\alpha}.
\tag{27}
\]
Insert (27) into (26) to obtain
\[
\E\!\bigl[f(\tilde{x}^{s+1})-f(x^{*})\bigr]
\le \underbrace{\Bigl[\frac{1-(1-\alpha)^{m}}{m\alpha}
+ \frac{L^{2}\eta^{2}}{\alpha}\Bigl(1-\frac{1-(1-\alpha)^{m}}{m\alpha}\Bigr)\Bigr]}_{\displaystyle \rho}\;
\bigl(f(\tilde{x}^{s})-f(x^{*})\bigr).
\tag{28}
\]

\medskip
\noindent\textbf{Step 12 (Simplify $\rho$).}  
Using $0<\alpha\le 1$ and the inequality $1-(1-\alpha)^{m}\le m\alpha$, we bound
\[
\frac{1-(1-\alpha)^{m}}{m\alpha}\le 1,
\qquad
1-\frac{1-(1-\alpha)^{m}}{m\alpha}\le 1-\frac{1-(1-\alpha)^{m}}{m\alpha}\le 1.
\]
Hence a convenient upper bound is
\[
\rho\;\le\; \frac{1}{m} + \Bigl(1-\frac{1}{m}\Bigr)\bigl(1-2\mu\eta\bigr),
\tag{29}
\]
where we have used $\alpha = \eta\mu(1- L\eta)-L^{2}\eta^{2}\ge \mu\eta(1-2L\eta)$
(because $0<\eta<\frac{1}{2L}$).  This yields exactly the expression (5) stated in the theorem.

\medskip
\noindent\textbf{Step 13 (Choosing hyperparameters).}  
Take $\eta=\frac{1}{4L}$; then $1-2L\eta=\frac12$.
If additionally $m\ge \frac{2L}{\mu}$, the bound (29) becomes
\[
\rho\le \frac{1}{m}+ \Bigl(1-\frac{1}{m}\Bigr)\frac12
\;=\;\frac12+\frac{1}{2m}\;<\;1.
\]
Thus $\rho\in(0,1)$ and (6) follows by induction over $s$,
which completes the proof. \hfill $\square$

\section*{Rate Derivation (With Constants)}
From (7) we have the explicit linear rate
\[
\E\!\bigl[f(\tilde{x}^{S})-f(x^{*})\bigr]\;\le\;
\Bigl(\tfrac12+\tfrac{1}{2m}\Bigr)^{S}
\bigl[f(\tilde{x}^{0})-f(x^{*})\bigr],
\]
so after $S$ epochs the optimality gap decays as $\mathcal{O}\!\bigl((1-\Theta(1))^{S}\bigr)$.
Since each epoch costs $n+m$ gradient evaluations, the total
complexity to reach $\varepsilon$ accuracy is
\[
\mathcal{O}\!\bigl((n+m)\log(1/\varepsilon)\bigr).
\]

\section*{Special Cases}
\begin{enumerate}[label=\alph*.]
\item \textbf{Full‑gradient case ($m=0$).} The algorithm reduces to gradient descent
with rate $(1- \mu\eta)^{t}$.
\item \textbf{Single‑sample case ($n=1$).} The variance term vanishes; SVRG
becomes ordinary gradient descent and the bound simplifies to the classic
$\bigl(1-\mu\eta\bigr)^{t}$ rate.
\item \textbf{Large $m$ (e.g., $m=n$).} The contraction factor approaches
$1-2\mu\eta$, matching the optimal rate of the proximal point method.
\end{enumerate}

\end{document}